\section{The Transformer Family}\label{the-transformer-family}

Date: April 7, 2020 \textbar{} Estimated Reading Time: 25 min \textbar{} Author: Lilian Weng

{Table of Contents}

\begin{itemize}
\tightlist
\item
  \hyperref[notations]{Notations}
\end{itemize}

\hyperref[attention-and-self-attention]{Attention and Self-Attention}

\hyperref[multi-head-self-attention]{Multi-Head Self-Attention}

\hyperref[transformer]{Transformer}

\hyperref[adaptive-computation-time-act]{Adaptive Computation Time (ACT)}

\hyperref[improved-attention-span]{Improved Attention Span}

\begin{itemize}
\tightlist
\item
  \hyperref[longer-attention-span-transformer-xl]{Longer Attention Span (Transformer-XL)}
\item
  \hyperref[adaptive-attention-span]{Adaptive Attention Span}
\item
  \hyperref[localized-attention-span-image-transformer]{Localized Attention Span (Image Transformer)}
\end{itemize}

\hyperref[less-time-and-memory-cost]{Less Time and Memory Cost}

\begin{itemize}
\tightlist
\item
  \hyperref[sparse-attention-matrix-factorization-sparse-transformers]{Sparse Attention Matrix Factorization (Sparse Transformers)}
\item
  \hyperref[locality-sensitive-hashing-reformer]{Locality-Sensitive Hashing (Reformer)}
\end{itemize}

\hyperref[make-it-recurrent-universal-transformer]{Make it Recurrent (Universal Transformer)}

\hyperref[stabilization-for-rl-gtrxl]{Stabilization for RL (GTrXL)}

\hyperref[citation]{Citation}

\hyperref[reference]{Reference}

{{[}Updated on \hl{\textbf{2023-01-27}}: After almost three years, I did a big refactoring update of this post to incorporate a bunch of new Transformer models since 2020. The enhanced version of this post is here: \href{https://lilianweng.github.io/posts/2023-01-27-the-transformer-family-v2/}{\hl{\textbf{The Transformer Family Version 2.0}}}. Please refer to that post on this topic.{]}}\\

It has been almost two years since my last post on \href{https://lilianweng.github.io/posts/2018-06-24-attention/}{attention}. Recent progress on new and enhanced versions of Transformer motivates me to write another post on this specific topic, focusing on how the vanilla Transformer can be improved for longer-term attention span, less memory and computation consumption, RL task solving and more.

\subsection{\texorpdfstring{Notations\hyperref[notations]{\#}}{Notations\#}}\label{notations}

\begin{longtable}[]{@{}ll@{}}
\toprule\noalign{}
Symbol & Meaning \\
\midrule\noalign{}
\endhead
\bottomrule\noalign{}
\endlastfoot
\$d\$ & The model size / hidden state dimension / positional encoding size. \\
\$h\$ & The number of heads in multi-head attention layer. \\
\$L\$ & The segment length of input sequence. \\
\$\textbackslash mathbf\{X\} \textbackslash in \textbackslash mathbb\{R\}\^{}\{L \textbackslash times d\}\$ & The input sequence where each element has been mapped into an embedding vector of shape \$d\$, same as the model size. \\
\$\textbackslash mathbf\{W\}\^{}k \textbackslash in \textbackslash mathbb\{R\}\^{}\{d \textbackslash times d\_k\}\$ & The key weight matrix. \\
\$\textbackslash mathbf\{W\}\^{}q \textbackslash in \textbackslash mathbb\{R\}\^{}\{d \textbackslash times d\_k\}\$ & The query weight matrix. \\
\$\textbackslash mathbf\{W\}\^{}v \textbackslash in \textbackslash mathbb\{R\}\^{}\{d \textbackslash times d\_v\}\$ & The value weight matrix. Often we have \$d\_k = d\_v = d\$. \\
\$\textbackslash mathbf\{W\}\^{}k\_i, \textbackslash mathbf\{W\}\^{}q\_i \textbackslash in \textbackslash mathbb\{R\}\^{}\{d \textbackslash times d\_k/h\}; \textbackslash mathbf\{W\}\^{}v\_i \textbackslash in \textbackslash mathbb\{R\}\^{}\{d \textbackslash times d\_v/h\}\$ & The weight matrices per head. \\
\$\textbackslash mathbf\{W\}\^{}o \textbackslash in \textbackslash mathbb\{R\}\^{}\{d\_v \textbackslash times d\}\$ & The output weight matrix. \\
\$\textbackslash mathbf\{Q\} = \textbackslash mathbf\{X\}\textbackslash mathbf\{W\}\^{}q \textbackslash in \textbackslash mathbb\{R\}\^{}\{L \textbackslash times d\_k\}\$ & The query embedding inputs. \\
\$\textbackslash mathbf\{K\} = \textbackslash mathbf\{X\}\textbackslash mathbf\{W\}\^{}k \textbackslash in \textbackslash mathbb\{R\}\^{}\{L \textbackslash times d\_k\}\$ & The key embedding inputs. \\
\$\textbackslash mathbf\{V\} = \textbackslash mathbf\{X\}\textbackslash mathbf\{W\}\^{}v \textbackslash in \textbackslash mathbb\{R\}\^{}\{L \textbackslash times d\_v\}\$ & The value embedding inputs. \\
\$S\_i\$ & A collection of key positions for the \$i\$-th query \$\textbackslash mathbf\{q\}\_i\$ to attend to. \\
\$\textbackslash mathbf\{A\} \textbackslash in \textbackslash mathbb\{R\}\^{}\{L \textbackslash times L\}\$ & The self-attention matrix between a input sequence of length \$L\$ and itself. \$\textbackslash mathbf\{A\} = \textbackslash text\{softmax\}(\textbackslash mathbf\{Q\}\textbackslash mathbf\{K\}\^{}\textbackslash top / \textbackslash sqrt\{d\_k\})\$. \\
\$a\_\{ij\} \textbackslash in \textbackslash mathbf\{A\}\$ & The scalar attention score between query \$\textbackslash mathbf\{q\}\_i\$ and key \$\textbackslash mathbf\{k\}\_j\$. \\
\$\textbackslash mathbf\{P\} \textbackslash in \textbackslash mathbb\{R\}\^{}\{L \textbackslash times d\}\$ & position encoding matrix, where the \$i\$-th row \$\textbackslash mathbf\{p\}\_i\$ is the positional encoding for input \$\textbackslash mathbf\{x\}\_i\$. \\
\end{longtable}

\section{\texorpdfstring{Attention and Self-Attention\hyperref[attention-and-self-attention]{\#}}{Attention and Self-Attention\#}}\label{attention-and-self-attention}

\emph{Attention} is a mechanism in the neural network that a model can learn to make predictions by selectively attending to a given set of data. The amount of attention is quantified by learned weights and thus the output is usually formed as a weighted average.

\emph{Self-attention} is a type of attention mechanism where the model makes prediction for one part of a data sample using other parts of the observation about the same sample. Conceptually, it feels quite similar to \href{https://en.wikipedia.org/wiki/Non-local_means}{non-local means}. Also note that self-attention is permutation-invariant; in other words, it is an operation on sets.

There are various forms of attention / self-attention, Transformer (\href{https://arxiv.org/abs/1706.03762}{Vaswani et al., 2017}) relies on the \emph{scaled dot-product attention}: given a query matrix \$\textbackslash mathbf\{Q\}\$, a key matrix \$\textbackslash mathbf\{K\}\$ and a value matrix \$\textbackslash mathbf\{V\}\$, the output is a weighted sum of the value vectors, where the weight assigned to each value slot is determined by the dot-product of the query with the corresponding key:

\$\$ \textbackslash text\{Attention\}(\textbackslash mathbf\{Q\}, \textbackslash mathbf\{K\}, \textbackslash mathbf\{V\}) = \textbackslash text\{softmax\}(\textbackslash frac\{\textbackslash mathbf\{Q\} \{\textbackslash mathbf\{K\}\}\^{}\textbackslash top\}\{\textbackslash sqrt\{d\_k\}\})\textbackslash mathbf\{V\} \$\$

And for a query and a key vector \$\textbackslash mathbf\{q\}\_i, \textbackslash mathbf\{k\}\_j \textbackslash in \textbackslash mathbb\{R\}\^{}d\$ (row vectors in query and key matrices), we have a scalar score:

\$\$ a\_\{ij\} = \textbackslash text\{softmax\}(\textbackslash frac\{\textbackslash mathbf\{q\}\_i \{\textbackslash mathbf\{k\}\_j\}\^{}\textbackslash top\}\{\textbackslash sqrt\{d\_k\}\}) = \textbackslash frac\{\textbackslash exp(\textbackslash mathbf\{q\}\_i \{\textbackslash mathbf\{k\}\_j\}\^{}\textbackslash top)\}\{ \textbackslash sqrt\{d\_k\} \textbackslash sum\_\{r \textbackslash in S\_i\} \textbackslash exp(\textbackslash mathbf\{q\}\_i \{\textbackslash mathbf\{k\}\_r\}\^{}\textbackslash top) \} \$\$

See my old \href{https://lilianweng.github.io/posts/2018-06-24-attention/\#a-family-of-attention-mechanisms}{post} for other types of attention if interested.

\section{\texorpdfstring{Multi-Head Self-Attention\hyperref[multi-head-self-attention]{\#}}{Multi-Head Self-Attention\#}}\label{multi-head-self-attention}

The \emph{multi-head self-attention} module is a key component in Transformer. Rather than only computing the attention once, the multi-head mechanism splits the inputs into smaller chunks and then computes the scaled dot-product attention over each subspace in parallel. The independent attention outputs are simply concatenated and linearly transformed into expected dimensions.

\$\$ \textbackslash begin\{aligned\} \textbackslash text\{MultiHeadAttention\}(\textbackslash mathbf\{X\}\_q, \textbackslash mathbf\{X\}\_k, \textbackslash mathbf\{X\}\_v) \&= {[}\textbackslash text\{head\}\_1; \textbackslash dots; \textbackslash text\{head\}\_h{]} \textbackslash mathbf\{W\}\^{}o \textbackslash\textbackslash{} \textbackslash text\{where head\}\_i \&= \textbackslash text\{Attention\}(\textbackslash mathbf\{X\}\_q\textbackslash mathbf\{W\}\^{}q\_i, \textbackslash mathbf\{X\}\_k\textbackslash mathbf\{W\}\^{}k\_i, \textbackslash mathbf\{X\}\_v\textbackslash mathbf\{W\}\^{}v\_i) \textbackslash end\{aligned\} \$\$

where \${[}.;.{]}\$ is a concatenation operation. \$\textbackslash mathbf\{W\}\^{}q\_i, \textbackslash mathbf\{W\}\^{}k\_i \textbackslash in \textbackslash mathbb\{R\}\^{}\{d \textbackslash times d\_k/h\}, \textbackslash mathbf\{W\}\^{}v\_i \textbackslash in \textbackslash mathbb\{R\}\^{}\{d \textbackslash times d\_v/h\}\$ are weight matrices to map input embeddings of size \$L \textbackslash times d\$ into query, key and value matrices. And \$\textbackslash mathbf\{W\}\^{}o \textbackslash in \textbackslash mathbb\{R\}\^{}\{d\_v \textbackslash times d\}\$ is the output linear transformation. All the weights should be learned during training.

\begin{figure}
\centering
\includegraphics{multi-head-attention.png}
\caption{Illustration of the multi-head scaled dot-product attention mechanism. (Image source: Figure 2 in \href{https://arxiv.org/abs/1706.03762}{Vaswani, et al., 2017})}
\end{figure}

\section{\texorpdfstring{Transformer\hyperref[transformer]{\#}}{Transformer\#}}\label{transformer}

The \textbf{Transformer} (which will be referred to as ``vanilla Transformer'' to distinguish it from other enhanced versions; \href{https://arxiv.org/abs/1706.03762}{Vaswani, et al., 2017}) model has an encoder-decoder architecture, as commonly used in many \href{https://lilianweng.github.io/posts/2018-06-24-attention/\#born-for-translation}{NMT} models. Later simplified Transformer was shown to achieve great performance in language modeling tasks, like in encoder-only \href{https://lilianweng.github.io/posts/2019-01-31-lm/\#bert}{BERT} or decoder-only \href{https://lilianweng.github.io/posts/2019-01-31-lm/\#openai-gpt}{GPT}.

\textbf{Encoder-Decoder Architecture}

The \textbf{encoder} generates an attention-based representation with capability to locate a specific piece of information from a large context. It consists of a stack of 6 identity modules, each containing two submodules, a \emph{multi-head self-attention} layer and a \emph{point-wise} fully connected feed-forward network. By point-wise, it means that it applies the same linear transformation (with same weights) to each element in the sequence. This can also be viewed as a convolutional layer with filter size 1. Each submodule has a residual connection and layer normalization. All the submodules output data of the same dimension \$d\$.

The function of Transformer \textbf{decoder} is to retrieve information from the encoded representation. The architecture is quite similar to the encoder, except that the decoder contains two multi-head attention submodules instead of one in each identical repeating module. The first multi-head attention submodule is \emph{masked} to prevent positions from attending to the future.

\begin{figure}
\centering
\includegraphics{transformer.png}
\caption{The architecture of the vanilla Transformer model. (Image source: \href{https://lilianweng.github.io/posts/2018-06-24-attention/\#full-architecture}{Figure 17})}
\end{figure}

\textbf{Positional Encoding}

Because self-attention operation is permutation invariant, it is important to use proper \textbf{positional encoding}to provide \emph{order information} to the model. The positional encoding \$\textbackslash mathbf\{P\} \textbackslash in \textbackslash mathbb\{R\}\^{}\{L \textbackslash times d\}\$ has the same dimension as the input embedding, so it can be added on the input directly. The vanilla Transformer considered two types of encodings:

(1) \emph{Sinusoidal positional encoding} is defined as follows, given the token position \$i=1,\textbackslash dots,L\$ and the dimension \$\textbackslash delta=1,\textbackslash dots,d\$:

\$\$ \textbackslash text\{PE\}(i,\textbackslash delta) = \textbackslash begin\{cases\} \textbackslash sin(\textbackslash frac\{i\}\{10000\^{}\{2\textbackslash delta\textquotesingle/d\}\}) \& \textbackslash text\{if \} \textbackslash delta = 2\textbackslash delta\textquotesingle\textbackslash\textbackslash{} \textbackslash cos(\textbackslash frac\{i\}\{10000\^{}\{2\textbackslash delta\textquotesingle/d\}\}) \& \textbackslash text\{if \} \textbackslash delta = 2\textbackslash delta\textquotesingle{} + 1\textbackslash\textbackslash{} \textbackslash end\{cases\} \$\$

In this way each dimension of the positional encoding corresponds to a sinusoid of different wavelengths in different dimensions, from \$2\textbackslash pi\$ to \$10000 \textbackslash cdot 2\textbackslash pi\$.

\begin{figure}
\centering
\includegraphics{sinoidual-positional-encoding.png}
\caption{Sinusoidal positional encoding with \$L=32\$ and \$d=128\$. The value is between -1 (black) and 1 (white) and the value 0 is in gray.}
\end{figure}

(2) \emph{Learned positional encoding}, as its name suggested, assigns each element with a learned column vector which encodes its \emph{absolute} position (\href{https://arxiv.org/abs/1705.03122}{Gehring, et al. 2017}).

\textbf{Quick Follow-ups}

Following the vanilla Transformer, \href{https://arxiv.org/abs/1808.04444}{Al-Rfou et al. (2018)} added a set of auxiliary losses to enable training a deep Transformer model on character-level language modeling which outperformed LSTMs. Several types of auxiliary tasks are used:

\begin{itemize}
\tightlist
\item
  Instead of producing only one prediction at the sequence end, every \emph{immediate position} is also asked to make a correct prediction, forcing the model to predict given smaller contexts (e.g. first couple tokens at the beginning of a context window).
\item
  Each intermediate Transformer layer is used for making predictions as well. Lower layers are weighted to contribute less and less to the total loss as training progresses.
\item
  Each position in the sequence can predict multiple targets, i.e. two or more predictions of the future tokens.
\end{itemize}

\begin{figure}
\centering
\includegraphics{transformer-aux-losses.png}
\caption{Auxiliary prediction tasks used in deep Transformer for character-level language modeling. (Image source: \href{https://arxiv.org/abs/1808.04444}{Al-Rfou et al. (2018)})}
\end{figure}

\section{\texorpdfstring{Adaptive Computation Time (ACT)\hyperref[adaptive-computation-time-act]{\#}}{Adaptive Computation Time (ACT)\#}}\label{adaptive-computation-time-act}

\textbf{Adaptive Computation Time} (short for \textbf{ACT}; \href{https://arxiv.org/abs/1603.08983}{Graves, 2016}) is a mechanism for dynamically deciding how many computational steps are needed in a recurrent neural network. Here is a cool \href{https://distill.pub/2016/augmented-rnns/\#adaptive-computation-time}{tutorial} on ACT from distill.pub.

Let's say, we have a RNN model \$\textbackslash mathcal\{R\}\$ composed of input weights \$W\_x\$, a parametric state transition function \$\textbackslash mathcal\{S\}(.)\$, a set of output weights \$W\_y\$ and an output bias \$b\_y\$. Given an input sequence \$(x\_1, \textbackslash dots, x\_L)\$, the output sequence \$(y\_1, \textbackslash dots, y\_L)\$ is computed by:

\$\$ s\_t = \textbackslash mathcal\{S\}(s\_\{t-1\}, W\_x x\_t), \textbackslash quad y\_t = W\_y s\_t + b\_y\textbackslash quad\textbackslash text\{for \}t=1, \textbackslash dots, L \$\$

ACT enables the above RNN setup to perform a variable number of steps at each input element. Multiple computational steps lead to a sequence of intermediate states \$(s\_t\^{}1, \textbackslash dots, s\_t\^{}\{N(t)\})\$ and outputs \$(y\_t\^{}1, \textbackslash dots, y\_t\^{}\{N(t)\})\$ --- they all share the same state transition function \$\textbackslash mathcal\{S\}(.)\$, as well as the same output weights \$W\_y\$ and bias \$b\_y\$:

\$\$ \textbackslash begin\{aligned\} s\_t\^{}0 \&= s\_\{t-1\} \textbackslash\textbackslash{} s\_t\^{}n \&= \textbackslash mathcal\{S\}(s\_\{t\}\^{}\{n-1\}, x\_t\^{}n) = \textbackslash mathcal\{S\}(s\_\{t\}\^{}\{n-1\}, x\_t + \textbackslash delta\_\{n,1\}) \textbackslash text\{ for \} n=1, \textbackslash dots, N(t)\textbackslash\textbackslash{} y\_t\^{}n \&= W\_y s\_t\^{}n + b\_y \textbackslash end\{aligned\} \$\$

where \$\textbackslash delta\_\{n,1\}\$ is a binary flag indicating whether the input step has been incremented.

The number of steps \$N(t)\$ is determined by an extra sigmoidal halting unit \$h\$, with associated weight matrix \$W\_h\$ and bias \$b\_h\$, outputting a halting probability \$p\_t\^{}n\$ at immediate step \$n\$ for \$t\$-th input element:

\$\$ h\_t\^{}n = \textbackslash sigma(W\_h s\_t\^{}n + b\_h) \$\$

In order to allow the computation to halt after a single step, ACT introduces a small constant \$\textbackslash epsilon\$ (e.g. 0.01), so that whenever the cumulative probability goes above \$1-\textbackslash epsilon\$, the computation stops.

\$\$ \textbackslash begin\{aligned\} N(t) \&= \textbackslash min(\textbackslash min\textbackslash\{n\textquotesingle: \textbackslash sum\_\{n=1\}\^{}\{n\textquotesingle\} h\_t\^{}n \textbackslash geq 1 -\textbackslash epsilon\textbackslash\}, M) \textbackslash\textbackslash{} p\_t\^{}n \&= \textbackslash begin\{cases\} h\_t\^{}n \& \textbackslash text\{if \}n \textless{} N(t) \textbackslash\textbackslash{} R(t) = 1 - \textbackslash sum\_\{n=1\}\^{}\{N(t)-1\} h\_t\^{}n \& \textbackslash text\{if \}n= N(t)\textbackslash\textbackslash{} \textbackslash end\{cases\} \textbackslash end\{aligned\} \$\$

where \$M\$ is an upper limit for the number of immediate steps allowed.

The final state and output are mean-field updates:

\$\$ s\_t = \textbackslash sum\_\{n=1\}\^{}\{N(t)\} p\_t\^{}n s\_t\^{}n,\textbackslash quad y\_t = \textbackslash sum\_\{n=1\}\^{}\{N(t)\} p\_t\^{}n y\_t\^{}n \$\$

\begin{figure}
\centering
\includegraphics{ACT-computation-graph.png}
\caption{The computation graph of a RNN with ACT mechanism. (Image source: \href{https://arxiv.org/abs/1603.08983}{Graves, 2016})}
\end{figure}

To avoid unnecessary pondering over each input, ACT adds a \emph{ponder cost} \$\textbackslash mathcal\{P\}(x) = \textbackslash sum\_\{t=1\}\^{}L N(t) + R(t) \$ in the loss function to encourage a smaller number of intermediate computational steps.

\section{\texorpdfstring{Improved Attention Span\hyperref[improved-attention-span]{\#}}{Improved Attention Span\#}}\label{improved-attention-span}

The goal of improving attention span is to make the context that can be used in self-attention longer, more efficient and flexible.

\subsection{\texorpdfstring{Longer Attention Span (Transformer-XL)\hyperref[longer-attention-span-transformer-xl]{\#}}{Longer Attention Span (Transformer-XL)\#}}\label{longer-attention-span-transformer-xl}

The vanilla Transformer has a fixed and limited attention span. The model can only attend to other elements in the same segments during each update step and no information can flow across separated fixed-length segments.

This \emph{context segmentation} causes several issues:

\begin{itemize}
\tightlist
\item
  The model cannot capture very long term dependencies.
\item
  It is hard to predict the first few tokens in each segment given no or thin context.
\item
  The evaluation is expensive. Whenever the segment is shifted to the right by one, the new segment is re-processed from scratch, although there are a lot of overlapped tokens.
\end{itemize}

\textbf{Transformer-XL} (\href{https://arxiv.org/abs/1901.02860}{Dai et al., 2019}; ``XL'' means ``extra long'') solves the context segmentation problem with two main modifications:

\begin{enumerate}
\tightlist
\item
  Reusing hidden states between segments.
\item
  Adopting a new positional encoding that is suitable for reused states.
\end{enumerate}

\textbf{Hidden State Reuse}

The recurrent connection between segments is introduced into the model by continuously using the hidden states from the previous segments.

\begin{figure}
\centering
\includegraphics{transformer-XL-training.png}
\caption{A comparison between the training phrase of vanilla Transformer \& Transformer-XL with a segment length 4. (Image source: left part of Figure 2 in \href{https://arxiv.org/abs/1901.02860}{Dai et al., 2019}).}
\end{figure}

Let's label the hidden state of the \$n\$-th layer for the \$(\textbackslash tau + 1)\$-th segment in the model as \$\textbackslash mathbf\{h\}\_\{\textbackslash tau+1\}\^{}\{(n)\} \textbackslash in \textbackslash mathbb\{R\}\^{}\{L \textbackslash times d\}\$. In addition to the hidden state of the last layer for the same segment \$\textbackslash mathbf\{h\}\_\{\textbackslash tau+1\}\^{}\{(n-1)\}\$, it also depends on the hidden state of the same layer for the previous segment \$\textbackslash mathbf\{h\}\_\{\textbackslash tau\}\^{}\{(n)\}\$. By incorporating information from the previous hidden states, the model extends the attention span much longer in the past, over multiple segments.

\$\$ \textbackslash begin\{aligned\} \textbackslash color\{red\}\{\textbackslash widetilde\{\textbackslash mathbf\{h\}\}\_\{\textbackslash tau+1\}\^{}\{(n-1)\}\} \&= {[}\textbackslash text\{stop-gradient\}(\textbackslash mathbf\{h\}\_\{\textbackslash tau\}\^{}\{(n-1)\}) \textbackslash circ \textbackslash mathbf\{h\}\_\{\textbackslash tau+1\}\^{}\{(n-1)\}{]} \textbackslash\textbackslash{} \textbackslash mathbf\{Q\}\_\{\textbackslash tau+1\}\^{}\{(n)\} \&= \textbackslash mathbf\{h\}\_\{\textbackslash tau+1\}\^{}\{(n-1)\}\textbackslash mathbf\{W\}\^{}q \textbackslash\textbackslash{} \textbackslash mathbf\{K\}\_\{\textbackslash tau+1\}\^{}\{(n)\} \&= \textbackslash color\{red\}\{\textbackslash widetilde\{\textbackslash mathbf\{h\}\}\_\{\textbackslash tau+1\}\^{}\{(n-1)\}\} \textbackslash mathbf\{W\}\^{}k \textbackslash\textbackslash{} \textbackslash mathbf\{V\}\_\{\textbackslash tau+1\}\^{}\{(n)\} \&= \textbackslash color\{red\}\{\textbackslash widetilde\{\textbackslash mathbf\{h\}\}\_\{\textbackslash tau+1\}\^{}\{(n-1)\}\} \textbackslash mathbf\{W\}\^{}v \textbackslash\textbackslash{} \textbackslash mathbf\{h\}\_\{\textbackslash tau+1\}\^{}\{(n)\} \&= \textbackslash text\{transformer-layer\}(\textbackslash mathbf\{Q\}\_\{\textbackslash tau+1\}\^{}\{(n)\}, \textbackslash mathbf\{K\}\_\{\textbackslash tau+1\}\^{}\{(n)\}, \textbackslash mathbf\{V\}\_\{\textbackslash tau+1\}\^{}\{(n)\}) \textbackslash end\{aligned\} \$\$

Note that both key and value rely on the extended hidden state, while the query only consumes hidden state at current step. The concatenation operation \${[}. \textbackslash circ .{]}\$ is along the sequence length dimension.

\textbf{Relative Positional Encoding}

In order to work with this new form of attention span, Transformer-XL proposed a new type of positional encoding. If using the same approach by vanilla Transformer and encoding the absolute position, the previous and current segments will be assigned with the same encoding, which is undesired.

To keep the positional information flow coherently across segments, Transformer-XL encodes the \emph{relative} position instead, as it could be sufficient enough to know the position offset for making good predictions, i.e. \$i-j\$, between one key vector \$\textbackslash mathbf\{k\}\_\{\textbackslash tau, j\}\$ and its query \$\textbackslash mathbf\{q\}\_\{\textbackslash tau, i\}\$.

If omitting the scalar \$1/\textbackslash sqrt\{d\_k\}\$ and the normalizing term in softmax but including positional encodings, we can write the attention score between query at position \$i\$ and key at position \$j\$ as:

\$\$ \textbackslash begin\{aligned\} a\_\{ij\} \&= \textbackslash mathbf\{q\}\_i \{\textbackslash mathbf\{k\}\_j\}\^{}\textbackslash top = (\textbackslash mathbf\{x\}\_i + \textbackslash mathbf\{p\}\_i)\textbackslash mathbf\{W\}\^{}q ((\textbackslash mathbf\{x\}\_j + \textbackslash mathbf\{p\}\_j)\textbackslash mathbf\{W\}\^{}k)\^{}\textbackslash top \textbackslash\textbackslash{} \&= \textbackslash mathbf\{x\}\_i\textbackslash mathbf\{W\}\^{}q \{\textbackslash mathbf\{W\}\^{}k\}\^{}\textbackslash top\textbackslash mathbf\{x\}\_j\^{}\textbackslash top + \textbackslash mathbf\{x\}\_i\textbackslash mathbf\{W\}\^{}q \{\textbackslash mathbf\{W\}\^{}k\}\^{}\textbackslash top\textbackslash mathbf\{p\}\_j\^{}\textbackslash top + \textbackslash mathbf\{p\}\_i\textbackslash mathbf\{W\}\^{}q \{\textbackslash mathbf\{W\}\^{}k\}\^{}\textbackslash top\textbackslash mathbf\{x\}\_j\^{}\textbackslash top + \textbackslash mathbf\{p\}\_i\textbackslash mathbf\{W\}\^{}q \{\textbackslash mathbf\{W\}\^{}k\}\^{}\textbackslash top\textbackslash mathbf\{p\}\_j\^{}\textbackslash top \textbackslash end\{aligned\} \$\$

Transformer-XL reparameterizes the above four terms as follows:

\$\$ a\_\{ij\}\^{}\textbackslash text\{rel\} = \textbackslash underbrace\{ \textbackslash mathbf\{x\}\_i\textbackslash mathbf\{W\}\^{}q \textbackslash color\{blue\}\{ \{\textbackslash mathbf\{W\}\_E\^{}k\}\^{}\textbackslash top \} \textbackslash mathbf\{x\}\_j\^{}\textbackslash top \}\_\textbackslash text\{content-based addressing\} + \textbackslash underbrace\{ \textbackslash mathbf\{x\}\_i\textbackslash mathbf\{W\}\^{}q \textbackslash color\{blue\}\{ \{\textbackslash mathbf\{W\}\_R\^{}k\}\^{}\textbackslash top \} \textbackslash color\{green\}\{\textbackslash mathbf\{r\}\_\{i-j\}\^{}\textbackslash top\} \}\_\textbackslash text\{content-dependent positional bias\} + \textbackslash underbrace\{ \textbackslash color\{red\}\{\textbackslash mathbf\{u\}\} \textbackslash color\{blue\}\{ \{\textbackslash mathbf\{W\}\_E\^{}k\}\^{}\textbackslash top \} \textbackslash mathbf\{x\}\_j\^{}\textbackslash top \}\_\textbackslash text\{global content bias\} + \textbackslash underbrace\{ \textbackslash color\{red\}\{\textbackslash mathbf\{v\}\} \textbackslash color\{blue\}\{ \{\textbackslash mathbf\{W\}\_R\^{}k\}\^{}\textbackslash top \} \textbackslash color\{green\}\{\textbackslash mathbf\{r\}\_\{i-j\}\^{}\textbackslash top\} \}\_\textbackslash text\{global positional bias\} \$\$

\begin{itemize}
\tightlist
\item
  Replace \$\textbackslash mathbf\{p\}\_j\$ with relative positional encoding \$\textbackslash mathbf\{r\}\_\{i-j\} \textbackslash in \textbackslash mathbf\{R\}\^{}\{d\}\$;
\item
  Replace \$\textbackslash mathbf\{p\}\_i\textbackslash mathbf\{W\}\^{}q\$ with two trainable parameters \$\textbackslash mathbf\{u\}\$ (for content) and \$\textbackslash mathbf\{v\}\$ (for location) in two different terms;
\item
  Split \$\textbackslash mathbf\{W\}\^{}k\$ into two matrices, \$\textbackslash mathbf\{W\}\^{}k\_E\$ for content information and \$\textbackslash mathbf\{W\}\^{}k\_R\$ for location information.
\end{itemize}

\subsection{\texorpdfstring{Adaptive Attention Span\hyperref[adaptive-attention-span]{\#}}{Adaptive Attention Span\#}}\label{adaptive-attention-span}

One key advantage of Transformer is the capability of capturing long-term dependencies. Depending on the context, the model may prefer to attend further sometime than others; or one attention head may had different attention pattern from the other. If the attention span could adapt its length flexibly and only attend further back when needed, it would help reduce both computation and memory cost to support longer maximum context size in the model.

This is the motivation for \textbf{Adaptive Attention Span}. \href{https://arxiv.org/abs/1905.07799}{Sukhbaatar, et al., (2019)} proposed a self-attention mechanism that seeks an optimal attention span. They hypothesized that different attention heads might assign scores differently within the same context window (See Fig. 7) and thus the optimal span would be trained separately per head.

\begin{figure}
\centering
\includegraphics{attention-per-head.png}
\caption{Two attention heads in the same model, A \& B, assign attention differently within the same context window. Head A attends more to the recent tokens, while head B look further back into the past uniformly. (Image source: \href{https://arxiv.org/abs/1905.07799}{Sukhbaatar, et al. 2019})}
\end{figure}

Given the \$i\$-th token, we need to compute the attention weights between this token and other keys at positions \$j \textbackslash in S\_i\$, where \$S\_i\$ defineds the \$i\$-th token's context window.

\$\$ \textbackslash begin\{aligned\} e\_\{ij\} \&= \textbackslash mathbf\{q\}\_i \{\textbackslash mathbf\{k\}\_j\}\^{}\textbackslash top \textbackslash\textbackslash{} a\_\{ij\} \&= \textbackslash text\{softmax\}(e\_\{ij\}) = \textbackslash frac\{\textbackslash exp(e\_\{ij\})\}\{\textbackslash sum\_\{r=i-s\}\^{}\{i-1\} \textbackslash exp(e\_\{ir\})\} \textbackslash\textbackslash{} \textbackslash mathbf\{y\}\_i \&= \textbackslash sum\_\{r=i-s\}\^{}\{i-1\}a\_\{ir\}\textbackslash mathbf\{v\}\_r = \textbackslash sum\_\{r=i-s\}\^{}\{i-1\}a\_\{ir\}\textbackslash mathbf\{x\}\_r\textbackslash mathbf\{W\}\^{}v \textbackslash end\{aligned\} \$\$

A \emph{soft mask function} \$m\_z\$ is added to control for an effective adjustable attention span, which maps the distance between query and key into a {[}0, 1{]} value. \$m\_z\$ is parameterized by \$z \textbackslash in {[}0, s{]}\$ and \$z\$ is to be learned:

\$\$ m\_z(x) = \textbackslash text\{clamp\}(\textbackslash frac\{1\}\{R\}(R+z-x), 0, 1) \$\$

where \$R\$ is a hyper-parameter which defines the softness of \$m\_z\$.

\begin{figure}
\centering
\includegraphics{soft-masking-function.png}
\caption{The soft masking function used in the adaptive attention span. (Image source: \href{https://arxiv.org/abs/1905.07799}{Sukhbaatar, et al. 2019}.)}
\end{figure}

The soft mask function is applied to the softmax elements in the attention weights:

\$\$ a\_\{ij\} = \textbackslash frac\{m\_z(i-j)\textbackslash exp(s\_\{ij\})\}\{\textbackslash sum\_\{r=i-s\}\^{}\{i-1\}m\_z(i-r) \textbackslash exp(s\_\{ir\})\} \$\$

In the above equation, \$z\$ is differentiable so it is trained jointly with other parts of the model. Parameters \$z\^{}\{(i)\}, i=1, \textbackslash dots, h\$ are learned \emph{separately per head}. Moreover, the loss function has an extra L1 penalty on \$\textbackslash sum\_\{i=1\}\^{}h z\^{}\{(i)\}\$.

Using \hyperref[adaptive-computation-time-act]{Adaptive Computation Time}, the approach can be further enhanced to have flexible attention span length, adaptive to the current input dynamically. The span parameter \$z\_t\$ of an attention head at time \$t\$ is a sigmoidal function, \$z\_t = S \textbackslash sigma(\textbackslash mathbf\{v\} \textbackslash cdot \textbackslash mathbf\{x\}\_t +b)\$, where the vector \$\textbackslash mathbf\{v\}\$ and the bias scalar \$b\$ are learned jointly with other parameters.

In the experiments of Transformer with adaptive attention span, \href{https://arxiv.org/abs/1905.07799}{Sukhbaatar, et al. (2019)} found a general tendency that lower layers do not require very long attention spans, while a few attention heads in higher layers may use exceptionally long spans. Adaptive attention span also helps greatly reduce the number of FLOPS, especially in a big model with many attention layers and a large context length.

\subsection{\texorpdfstring{Localized Attention Span (Image Transformer)\hyperref[localized-attention-span-image-transformer]{\#}}{Localized Attention Span (Image Transformer)\#}}\label{localized-attention-span-image-transformer}

The original, also the most popular, use case for Transformer is to do language modeling. The text sequence is one-dimensional in a clearly defined chronological order and thus the attention span grows linearly with increased context size.

However, if we want to use Transformer on images, it is unclear how to define the scope of context or the order. \textbf{Image Transformer} (\href{https://arxiv.org/abs/1802.05751}{Parmer, et al 2018}) embraces a formulation of image generation similar to sequence modeling within the Transformer framework. Additionally, Image Transformer restricts the self-attention span to only \emph{local} neighborhoods, so that the model can scale up to process more images in parallel and keep the likelihood loss tractable.

The encoder-decoder architecture remains for image-conditioned generation:

\begin{itemize}
\tightlist
\item
  The encoder generates a contextualized, per-pixel-channel representation of the source image;
\item
  The decoder \emph{autoregressively} generates an output image, one channel per pixel at each time step.
\end{itemize}

Let's label the representation of the current pixel to be generated as the query \$\textbackslash mathbf\{q\}\$. Other positions whose representations will be used for computing \$\textbackslash mathbf\{q\}\$ are key vector \$\textbackslash mathbf\{k\}\_1, \textbackslash mathbf\{k\}\_2, \textbackslash dots\$ and they together form a memory matrix \$\textbackslash mathbf\{M\}\$. The scope of \$\textbackslash mathbf\{M\}\$ defines the context window for pixel query \$\textbackslash mathbf\{q\}\$.

Image Transformer introduced two types of localized \$\textbackslash mathbf\{M\}\$, as illustrated below.

\begin{figure}
\centering
\includegraphics{image-transformer-attention.png}
\caption{Illustration of 1D and 2D attention span for visual inputs in Image Transformer. The black line marks a query block and the cyan outlines the actual attention span for pixel q. (Image source: Figure 2 in \href{https://arxiv.org/abs/1802.05751}{Parmer et al, 2018})}
\end{figure}

(1) \emph{1D Local Attention}: The input image is flattened in the \href{https://en.wikipedia.org/wiki/Raster_scan\#Scanning_pattern}{raster scanning} order, that is, from left to right and top to bottom. The linearized image is then partitioned into non-overlapping query blocks. The context window consists of pixels in the same query block as \$\textbackslash mathbf\{q\}\$ and a fixed number of additional pixels generated before this query block.

(2) \emph{2D Local Attention}: The image is partitioned into multiple non-overlapping rectangular query blocks. The query pixel can attend to all others in the same memory blocks. To make sure the pixel at the top-left corner can also have a valid context window, the memory block is extended to the top, left and right by a fixed amount, respectively.

\section{\texorpdfstring{Less Time and Memory Cost\hyperref[less-time-and-memory-cost]{\#}}{Less Time and Memory Cost\#}}\label{less-time-and-memory-cost}

This section introduces several improvements made on Transformer to reduce the computation time and memory consumption.

\subsection{\texorpdfstring{Sparse Attention Matrix Factorization (Sparse Transformers)\hyperref[sparse-attention-matrix-factorization-sparse-transformers]{\#}}{Sparse Attention Matrix Factorization (Sparse Transformers)\#}}\label{sparse-attention-matrix-factorization-sparse-transformers}

The compute and memory cost of the vanilla Transformer grows quadratically with sequence length and thus it is hard to be applied on very long sequences.

\textbf{Sparse Transformer} (\href{https://arxiv.org/abs/1904.10509}{Child et al., 2019}) introduced \emph{factorized self-attention}, through sparse matrix factorization, making it possible to train dense attention networks with hundreds of layers on sequence length up to 16,384, which would be infeasible on modern hardware otherwise.

Given a set of attention connectivity pattern \$\textbackslash mathcal\{S\} = \textbackslash\{S\_1, \textbackslash dots, S\_n\textbackslash\}\$, where each \$S\_i\$ records a set of key positions that the \$i\$-th query vector attends to.

\$\$ \textbackslash begin\{aligned\} \textbackslash text\{Attend\}(\textbackslash mathbf\{X\}, \textbackslash mathcal\{S\}) \&= \textbackslash Big( a(\textbackslash mathbf\{x\}\_i, S\_i) \textbackslash Big)\_\{i \textbackslash in \textbackslash\{1, \textbackslash dots, L\textbackslash\}\} \textbackslash\textbackslash{} \textbackslash text\{ where \} a(\textbackslash mathbf\{x\}\_i, S\_i) \&= \textbackslash text\{softmax\}\textbackslash Big(\textbackslash frac\{(\textbackslash mathbf\{x\}\_i \textbackslash mathbf\{W\}\^{}q)(\textbackslash mathbf\{x\}\_j \textbackslash mathbf\{W\}\^{}k)\_\{j \textbackslash in S\_i\}\^{}\textbackslash top\}\{\textbackslash sqrt\{d\_k\}\}\textbackslash Big) (\textbackslash mathbf\{x\}\_j \textbackslash mathbf\{W\}\^{}v)\_\{j \textbackslash in S\_i\} \textbackslash end\{aligned\} \$\$

Note that although the size of \$S\_i\$ is not fixed, \$a(\textbackslash mathbf\{x\}\_i, S\_i)\$ is always of size \$d\_v\$ and thus \$\textbackslash text\{Attend\}(\textbackslash mathbf\{X\}, \textbackslash mathcal\{S\}) \textbackslash in \textbackslash mathbb\{R\}\^{}\{L \textbackslash times d\_v\}\$.

In anto-regressive models, one attention span is defined as \$S\_i = \textbackslash\{j: j \textbackslash leq i\textbackslash\}\$ as it allows each token to attend to all the positions in the past.

In factorized self-attention, the set \$S\_i\$ is decomposed into a \emph{tree} of dependencies, such that for every pair of \$(i, j)\$ where \$j \textbackslash leq i\$, there is a path connecting \$i\$ back to \$j\$ and \$i\$ can attend to \$j\$ either directly or indirectly.

Precisely, the set \$S\_i\$ is divided into \$p\$ \emph{non-overlapping} subsets, where the \$m\$-th subset is denoted as \$A\^{}\{(m)\}\_i \textbackslash subset S\_i, m = 1,\textbackslash dots, p\$. Therefore the path between the output position \$i\$ and any \$j\$ has a maximum length \$p + 1\$. For example, if \$(j, a, b, c, \textbackslash dots, i)\$ is a path of indices between \$i\$ and \$j\$, we would have \$j \textbackslash in A\_a\^{}\{(1)\}, a \textbackslash in A\_b\^{}\{(2)\}, b \textbackslash in A\_c\^{}\{(3)\}, \textbackslash dots\$, so on and so forth.

\textbf{Sparse Factorized Attention}

Sparse Transformer proposed two types of fractorized attention. It is easier to understand the concepts as illustrated in Fig. 10 with 2D image inputs as examples.

\begin{figure}
\centering
\includegraphics{sparse-attention.png}
\caption{The top row illustrates the attention connectivity patterns in (a) Transformer, (b) Sparse Transformer with strided attention, and (c) Sparse Transformer with fixed attention. The bottom row contains corresponding self-attention connectivity matrices. Note that the top and bottom rows are not in the same scale. (Image source: \href{https://arxiv.org/abs/1904.10509}{Child et al., 2019} + a few of extra annotations.)}
\end{figure}

(1) \emph{Strided} attention with stride \$\textbackslash ell \textbackslash sim \textbackslash sqrt\{n\}\$. This works well with image data as the structure is aligned with strides. In the image case, each pixel would attend to all the previous \$\textbackslash ell\$ pixels in the raster scanning order (naturally cover the entire width of the image) and then those pixels attend to others in the same column (defined by another attention connectivity subset).

\$\$ \textbackslash begin\{aligned\} A\_i\^{}\{(1)\} \&= \textbackslash\{ t, t+1, \textbackslash dots, i\textbackslash\} \textbackslash text\{, where \} t = \textbackslash max(0, i - \textbackslash ell) \textbackslash\textbackslash{} A\_i\^{}\{(2)\} \&= \textbackslash\{j: (i-j) \textbackslash mod \textbackslash ell = 0\textbackslash\} \textbackslash end\{aligned\} \$\$

(2) \emph{Fixed} attention. A small set of tokens summarize previous locations and propagate that information to all future locations.

\$\$ \textbackslash begin\{aligned\} A\_i\^{}\{(1)\} \&= \textbackslash\{j: \textbackslash lfloor \textbackslash frac\{j\}\{\textbackslash ell\} \textbackslash rfloor = \textbackslash lfloor \textbackslash frac\{i\}\{\textbackslash ell\} \textbackslash rfloor \textbackslash\} \textbackslash\textbackslash{} A\_i\^{}\{(2)\} \&= \textbackslash\{j: j \textbackslash mod \textbackslash ell \textbackslash in \textbackslash\{\textbackslash ell-c, \textbackslash dots, \textbackslash ell-1\textbackslash\} \textbackslash\} \textbackslash end\{aligned\} \$\$

where \$c\$ is a hyperparameter. If \$c=1\$, it restricts the representation whereas many depend on a few positions. The paper chose \$c\textbackslash in \textbackslash\{ 8, 16, 32 \textbackslash\}\$ for \$\textbackslash ell \textbackslash in \textbackslash\{ 128, 256 \textbackslash\}\$.

\textbf{Use Factorized Self-Attention in Transformer}

There are three ways to use sparse factorized attention patterns in Transformer architecture:

\begin{enumerate}
\tightlist
\item
  One attention type per residual block and then interleave them,\\
  \$\textbackslash text\{attention\}(\textbackslash mathbf\{X\}) = \textbackslash text\{Attend\}(\textbackslash mathbf\{X\}, A\^{}\{(n \textbackslash mod p)\}) \textbackslash mathbf\{W\}\^{}o\$, where \$n\$ is the index of the current residual block.
\item
  Set up a single head which attends to locations that all the factorized heads attend to,\\
  \$\textbackslash text\{attention\}(\textbackslash mathbf\{X\}) = \textbackslash text\{Attend\}(\textbackslash mathbf\{X\}, \textbackslash cup\_\{m=1\}\^{}p A\^{}\{(m)\}) \textbackslash mathbf\{W\}\^{}o \$.
\item
  Use a multi-head attention mechanism, but different from vanilla Transformer, each head might adopt a pattern presented above, 1 or 2. =\textgreater{} This option often performs the best.
\end{enumerate}

Sparse Transformer also proposed a set of changes so as to train the Transformer up to hundreds of layers, including gradient checkpointing, recomputing attention \& FF layers during the backward pass, mixed precision training, efficient block-sparse implementation, etc. Please check the \href{https://arxiv.org/abs/1904.10509}{paper} for more details.

\subsection{\texorpdfstring{Locality-Sensitive Hashing (Reformer)\hyperref[locality-sensitive-hashing-reformer]{\#}}{Locality-Sensitive Hashing (Reformer)\#}}\label{locality-sensitive-hashing-reformer}

The improvements proposed by the \textbf{Reformer} model (\href{https://arxiv.org/abs/2001.04451}{Kitaev, et al. 2020}) aim to solve the following pain points in Transformer:

\begin{itemize}
\tightlist
\item
  Memory in a model with \$N\$ layers is \$N\$-times larger than in a single-layer model because we need to store activations for back-propagation.
\item
  The intermediate FF layers are often quite large.
\item
  The attention matrix on sequences of length \$L\$ often requires \$O(L\^{}2)\$ in both memory and time.
\end{itemize}

Reformer proposed two main changes:

\begin{enumerate}
\tightlist
\item
  Replace the dot-product attention with \emph{locality-sensitive hashing (LSH) attention}, reducing the complexity from \$O(L\^{}2)\$ to \$O(L\textbackslash log L)\$.
\item
  Replace the standard residual blocks with \emph{reversible residual layers}, which allows storing activations only once during training instead of \$N\$ times (i.e. proportional to the number of layers).
\end{enumerate}

\phantomsection\label{LSH}{}\textbf{Locality-Sensitive Hashing Attention}

In \$\textbackslash mathbf\{Q\} \textbackslash mathbf\{K\}\^{}\textbackslash top\$ part of the \hyperref[attention-and-self-attention]{attention formula}, we are only interested in the largest elements as only large elements contribute a lot after softmax. For each query \$\textbackslash mathbf\{q\}\_i \textbackslash in \textbackslash mathbf\{Q\}\$, we are looking for row vectors in \$\textbackslash mathbf\{K\}\$ closest to \$\textbackslash mathbf\{q\}\_i\$. In order to find nearest neighbors quickly in high-dimensional space, Reformer incorporates \href{https://en.wikipedia.org/wiki/Locality-sensitive_hashing}{Locality-Sensitive Hashing (LSH)} into its attention mechanism.

A hashing scheme \$x \textbackslash mapsto h(x)\$ is \emph{locality-sensitive} if it preserves the distancing information between data points, such that close vectors obtain similar hashes while distant vectors have very different ones. The Reformer adopts a hashing scheme as such, given a fixed random matrix \$\textbackslash mathbf\{R\} \textbackslash in \textbackslash mathbb\{R\}\^{}\{d \textbackslash times b/2\}\$ (where \$b\$ is a hyperparam), the hash function is \$h(x) = \textbackslash arg\textbackslash max({[}xR; −xR{]})\$.

\begin{figure}
\centering
\includegraphics{LSH-attention-matrix.png}
\caption{Illustration of Locality-Sensitive Hashing (LSH) attention. (Image source: right part of Figure 1 in \href{https://arxiv.org/abs/2001.04451}{Kitaev, et al. 2020}).}
\end{figure}

In LSH attention, a query can only attend to positions in the same hashing bucket, \$S\_i = \textbackslash\{j: h(\textbackslash mathbf\{q\}\_i) = h(\textbackslash mathbf\{k\}\_j)\textbackslash\}\$. It is carried out in the following process, as illustrated in Fig. 11:

\begin{itemize}
\tightlist
\item
  (a) The attention matrix for full attention is often sparse.
\item
  (b) Using LSH, we can sort the keys and queries to be aligned according to their hash buckets.
\item
  (c) Set \$\textbackslash mathbf\{Q\} = \textbackslash mathbf\{K\}\$ (precisely \$\textbackslash mathbf\{k\}\_j = \textbackslash mathbf\{q\}\_j / \textbar\textbackslash mathbf\{q\}\_j\textbar\$), so that there are equal numbers of keys and queries in one bucket, easier for batching. Interestingly, this ``shared-QK'' config does not affect the performance of the Transformer.
\item
  (d) Apply batching where chunks of \$m\$ consecutive queries are grouped together.
\end{itemize}

\begin{figure}
\centering
\includegraphics{LSH-attention.png}
\caption{The LSH attention consists of 4 steps: bucketing, sorting, chunking, and attention computation. (Image source: left part of Figure 1 in \href{https://arxiv.org/abs/2001.04451}{Kitaev, et al. 2020}).}
\end{figure}

\textbf{Reversible Residual Network}

Another improvement by Reformer is to use \emph{reversible residual layers} (\href{https://arxiv.org/abs/1707.04585}{Gomez et al. 2017}). The motivation for reversible residual network is to design the architecture in a way that activations at any given layer can be recovered from the activations at the following layer, using only the model parameters. Hence, we can save memory by recomputing the activation during backprop rather than storing all the activations.

Given a layer \$x \textbackslash mapsto y\$, the normal residual layer does \$y = x + F(x)\$, but the reversible layer splits both input and output into pairs \$(x\_1, x\_2) \textbackslash mapsto (y\_1, y\_2)\$ and then executes the following:

\$\$ y\_1 = x\_1 + F(x\_2),\textbackslash; y\_2 = x\_2 + G(y\_1) \$\$

and reversing is easy:

\$\$ x\_2 = y\_2 - G(y\_1), \textbackslash; x\_1 = y\_1 − F(x\_2) \$\$

Reformer applies the same idea to Transformer by combination attention (\$F\$) and feed-forward layers (\$G\$) within a reversible net block:

\$\$ Y\_1 = X\_1 + \textbackslash text\{Attention\}(X\_2), \textbackslash; Y\_2 = X\_2 + \textbackslash text\{FeedForward\}(Y\_1) \$\$

The memory can be further reduced by chunking the feed-forward computation:

\$\$ Y\_2 = {[}Y\_2\^{}\{(1)\}; \textbackslash dots; Y\_2\^{}\{(c)\}{]} = {[}X\_2\^{}\{(1)\} + \textbackslash text\{FeedForward\}(Y\_1\^{}\{(1)\}); \textbackslash dots; X\_2\^{}\{(c)\} + \textbackslash text\{FeedForward\}(Y\_1\^{}\{(c)\}){]} \$\$

The resulting reversible Transformer does not need to store activation in every layer.

\section{\texorpdfstring{Make it Recurrent (Universal Transformer)\hyperref[make-it-recurrent-universal-transformer]{\#}}{Make it Recurrent (Universal Transformer)\#}}\label{make-it-recurrent-universal-transformer}

The \textbf{Universal Transformer} (\href{https://arxiv.org/abs/1807.03819}{Dehghani, et al. 2019}) combines self-attention in Transformer with the recurrent mechanism in RNN, aiming to benefit from both a long-term global receptive field of Transformer and learned inductive biases of RNN.

Rather than going through a fixed number of layers, Universal Transformer dynamically adjusts the number of steps using \hyperref[adaptive-computation-time-act]{adaptive computation time}. If we fix the number of steps, an Universal Transformer is equivalent to a multi-layer Transformer with shared parameters across layers.

On a high level, the universal transformer can be viewed as a recurrent function for learning the hidden state representation per token. The recurrent function evolves in parallel across token positions and the information between positions is shared through self-attention.

\begin{figure}
\centering
\includegraphics{universal-transformer-loop.png}
\caption{How the Universal Transformer refines a set of hidden state representations repeatedly for every position in parallel. (Image source: Figure 1 in \href{https://arxiv.org/abs/1807.03819}{Dehghani, et al. 2019}).}
\end{figure}

Given an input sequence of length \$L\$, Universal Transformer iteratively updates the representation \$\textbackslash mathbf\{H\}\^{}t \textbackslash in \textbackslash mathbb\{R\}\^{}\{L \textbackslash times d\}\$ at step \$t\$ for an adjustable number of steps. At step 0, \$\textbackslash mathbf\{H\}\^{}0\$ is initialized to be same as the input embedding matrix. All the positions are processed in parallel in the multi-head self-attention mechanism and then go through a recurrent transition function.

\$\$ \textbackslash begin\{aligned\} \textbackslash mathbf\{A\}\^{}t \&= \textbackslash text\{LayerNorm\}(\textbackslash mathbf\{H\}\^{}\{t-1\} + \textbackslash text\{MultiHeadAttention\}(\textbackslash mathbf\{H\}\^{}\{t-1\} + \textbackslash mathbf\{P\}\^{}t) \textbackslash\textbackslash{} \textbackslash mathbf\{H\}\^{}t \&= \textbackslash text\{LayerNorm\}(\textbackslash mathbf\{A\}\^{}\{t-1\} + \textbackslash text\{Transition\}(\textbackslash mathbf\{A\}\^{}t)) \textbackslash end\{aligned\} \$\$

where \$\textbackslash text\{Transition\}(.)\$ is either a \href{https://arxiv.org/abs/1610.02357}{separable convolution} or a fully-connected neural network that consists of two position-wise (i.e. applied to each row of \$\textbackslash mathbf\{A\}\^{}t\$ individually) affine transformation + one ReLU.

The positional encoding \$\textbackslash mathbf\{P\}\^{}t\$ uses sinusoidal position signal but with an additional time dimension:

\$\$ \textbackslash text\{PE\}(i, t, \textbackslash delta) = \textbackslash begin\{cases\} \textbackslash sin(\textbackslash frac\{i\}\{10000\^{}\{2\textbackslash delta\textquotesingle/d\}\}) \textbackslash oplus \textbackslash sin(\textbackslash frac\{t\}\{10000\^{}\{2\textbackslash delta\textquotesingle/d\}\}) \& \textbackslash text\{if \} \textbackslash delta = 2\textbackslash delta\textquotesingle\textbackslash\textbackslash{} \textbackslash cos(\textbackslash frac\{i\}\{10000\^{}\{2\textbackslash delta\textquotesingle/d\}\}) \textbackslash oplus \textbackslash cos(\textbackslash frac\{t\}\{10000\^{}\{2\textbackslash delta\textquotesingle/d\}\}) \& \textbackslash text\{if \} \textbackslash delta = 2\textbackslash delta\textquotesingle{} + 1\textbackslash\textbackslash{} \textbackslash end\{cases\} \$\$

\begin{figure}
\centering
\includegraphics{universal-transformer.png}
\caption{A simplified illustration of Universal Transformer. The encoder and decoder share the same basic recurrent structure. But the decoder also attends to final encoder representation \$\textbackslash mathbf\{H\}\^{}T\$. (Image source: Figure 2 in \href{https://arxiv.org/abs/1807.03819}{Dehghani, et al. 2019})}
\end{figure}

In the adaptive version of Universal Transformer, the number of recurrent steps \$T\$ is dynamically determined by \hyperref[adaptive-computation-time-act]{ACT}. Each position is equipped with a dynamic ACT halting mechanism. Once a per-token recurrent block halts, it stops taking more recurrent updates but simply copies the current value to the next step until all the blocks halt or until the model reaches a maximum step limit.

\section{\texorpdfstring{Stabilization for RL (GTrXL)\hyperref[stabilization-for-rl-gtrxl]{\#}}{Stabilization for RL (GTrXL)\#}}\label{stabilization-for-rl-gtrxl}

The self-attention mechanism avoids compressing the whole past into a fixed-size hidden state and does not suffer from vanishing or exploding gradients as much as RNNs. Reinforcement Learning tasks can for sure benefit from these traits. \emph{However}, it is quite difficult to train Transformer even in supervised learning, let alone in the RL context. It could be quite challenging to stabilize and train a LSTM agent by itself, after all.

The \textbf{Gated Transformer-XL} (\textbf{GTrXL}; \href{https://arxiv.org/abs/1910.06764}{Parisotto, et al. 2019}) is one attempt to use Transformer for RL. GTrXL succeeded in stabilizing training with two changes on top of \hyperref[longer-attention-span-transformer-xl]{Transformer-XL}:

\begin{enumerate}
\tightlist
\item
  The layer normalization is only applied on the input stream in a residual module, but NOT on the shortcut stream. A key benefit to this reordering is to allow the original input to flow from the first to last layer.
\item
  The residual connection is replaced with a GRU-style (Gated Recurrent Unit; \href{https://arxiv.org/abs/1412.3555}{Chung et al., 2014}) \emph{gating} mechanism.
\end{enumerate}

\$\$ \textbackslash begin\{aligned\} r \&= \textbackslash sigma(W\_r\^{}\{(l)\} y + U\_r\^{}\{(l)\} x) \textbackslash\textbackslash{} z \&= \textbackslash sigma(W\_z\^{}\{(l)\} y + U\_z\^{}\{(l)\} x - b\_g\^{}\{(l)\}) \textbackslash\textbackslash{} \textbackslash hat\{h\} \&= \textbackslash tanh(W\_g\^{}\{(l)\} y + U\_g\^{}\{(l)\} (r \textbackslash odot x)) \textbackslash\textbackslash{} g\^{}\{(l)\}(x, y) \&= (1-z)\textbackslash odot x + z\textbackslash odot \textbackslash hat\{h\} \textbackslash end\{aligned\} \$\$

The gating function parameters are explicitly initialized to be close to an identity map - this is why there is a \$b\_g\$ term. A \$b\_g \textgreater{} 0\$ greatly helps with the learning speedup.

\begin{figure}
\centering
\includegraphics{gated-transformer-XL.png}
\caption{Comparison of the model architecture of Transformer-XL, Transformer-XL with the layer norm reordered, and Gated Transformer-XL. (Image source: Figure 1 in \href{https://arxiv.org/abs/1910.06764}{Parisotto, et al. 2019})}
\end{figure}

\section{\texorpdfstring{Citation\hyperref[citation]{\#}}{Citation\#}}\label{citation}

Cited as:

\begin{quote}
Weng, Lilian. (Apr 2020). The transformer family. Lil'Log. https://lilianweng.github.io/posts/2020-04-07-the-transformer-family/.
\end{quote}

Or

\begin{verbatim}
@article{weng2020transformer,
  title   = "The Transformer Family",
  author  = "Weng, Lilian",
  journal = "lilianweng.github.io",
  year    = "2020",
  month   = "Apr",
  url     = "https://lilianweng.github.io/posts/2020-04-07-the-transformer-family/"
}
\end{verbatim}

\section{\texorpdfstring{Reference\hyperref[reference]{\#}}{Reference\#}}\label{reference}

{[}1{]} Ashish Vaswani, et al. \href{http://papers.nips.cc/paper/7181-attention-is-all-you-need.pdf}{``Attention is all you need.''} NIPS 2017.

{[}2{]} Rami Al-Rfou, et al. \href{https://arxiv.org/abs/1808.04444}{``Character-level language modeling with deeper self-attention.''} AAAI 2019.

{[}3{]} Olah \& Carter, \href{http://doi.org/10.23915/disti}{``Attention and Augmented Recurrent Neural Networks''}, Distill, 2016.

{[}4{]} Sainbayar Sukhbaatar, et al. \href{https://arxiv.org/abs/1905.07799}{``Adaptive Attention Span in Transformers''}. ACL 2019.

{[}5{]} Rewon Child, et al. \href{https://arxiv.org/abs/1904.10509}{``Generating Long Sequences with Sparse Transformers''} arXiv:1904.10509 (2019).

{[}6{]} Nikita Kitaev, et al. \href{https://arxiv.org/abs/2001.04451}{``Reformer: The Efficient Transformer''} ICLR 2020.

{[}7{]} Alex Graves. (``Adaptive Computation Time for Recurrent Neural Networks''){[}https://arxiv.org/abs/1603.08983{]}

{[}8{]} Niki Parmar, et al. \href{https://arxiv.org/abs/1802.05751}{``Image Transformer''} ICML 2018.

{[}9{]} Zihang Dai, et al. \href{https://arxiv.org/abs/1901.02860}{``Transformer-XL: Attentive Language Models Beyond a Fixed-Length Context.''} ACL 2019.

{[}10{]} Aidan N. Gomez, et al. \href{https://arxiv.org/abs/1707.04585}{``The Reversible Residual Network: Backpropagation Without Storing Activations''} NIPS 2017.

{[}11{]} Mostafa Dehghani, et al. \href{https://arxiv.org/abs/1807.03819}{``Universal Transformers''} ICLR 2019.

{[}12{]} Emilio Parisotto, et al. \href{https://arxiv.org/abs/1910.06764}{``Stabilizing Transformers for Reinforcement Learning''} arXiv:1910.06764 (2019).

\begin{itemize}
\tightlist
\item
  \href{https://lilianweng.github.io/tags/architecture/}{Architecture}
\item
  \href{https://lilianweng.github.io/tags/attention/}{Attention}
\item
  \href{https://lilianweng.github.io/tags/transformer/}{Transformer}
\item
  \href{https://lilianweng.github.io/tags/foundation/}{Foundation}
\item
  \href{https://lilianweng.github.io/tags/reinforcement-learning/}{Reinforcement-Learning}
\end{itemize}

\href{https://lilianweng.github.io/posts/2020-06-07-exploration-drl/}{{« }\\
{Exploration Strategies in Deep Reinforcement Learning}} \href{https://lilianweng.github.io/posts/2020-01-29-curriculum-rl/}{{ »}\\
{Curriculum for Reinforcement Learning}}

